\chapter{Verification}
\label{ch:verif}

\section{Environment}

The target is behavioural RTL simulation with no vendor IP. The reference flow is
ModelSim; a second flow runs the same system testbench under Verilator to obtain
assertions and functional coverage, neither of which the free ModelSim edition
supports.

\begin{table}[H]
\centering
\caption{Verification collateral}
\label{tab:collateral}
\begin{tabular}{@{}lL{0.62\linewidth}@{}}
\toprule
\textbf{File} & \textbf{Role} \\
\midrule
\file{debug/hdl/tb\_cpu\_axi.v} & System bench: CPU, instruction and data memory models, two
  address decoders, the real PIC and machine timer, four protocol monitors. The bench plays the
  peripherals on \sig{irq\_src}. \\
\file{debug/hdl/tb\_pic.v} & Feature bench for the PIC alone, so each mechanism is a named check
  rather than something buried in a system run. \\
\file{debug/hdl/tb\_pic\_reset.v} & Reset bench for the PIC: every register's reset value, X-freedom
  including state with no register view, and an asynchronous reset taken from a fully configured,
  nested state. \\
\file{debug/hdl/tb\_pic\_ro.v} & Read-only and reserved-bit bench for the PIC: rejection and
  non-modification for every read-only object, and zero for every reserved position. \\
\file{debug/hdl/tb\_pic\_status.v} & \reg{SRCx\_STATUS} field by field, each in both directions,
  plus a per-clock-edge watch of the bump-escalation walk. \\
\file{debug/hdl/tb\_mtimer\_regs.v} & Machine-timer registers: reset values, exact counting rate,
  byte lanes, unmapped offsets, and the arm/fire/clear sequence. \\
\file{debug/hdl/tb\_csr\_ro.v} & CSR access classes kept distinct: read-only, WARL, tied off,
  absent. \\
\file{debug/hdl/tb\_traps.v} & One isolated directed test per trap cause, at CPU level, with
  \reg{mcause}, \reg{mepc}, \reg{mtval} and the handler entry address checked individually. \\
\file{debug/hdl/tb\_dual\_core.v} & Two cores sharing one memory behind a 2:1 arbiter. \\
\file{debug/hdl/axi\_lite\_monitor.v} & Passive protocol checker, instantiated on all four AXI links. \\
\file{debug/hdl/axi\_lite\_mem\_model.v} & Behavioural slave with configurable latency and seeded
  random \sig{READY} backpressure. \\
\file{debug/hdl/ck\_rst\_tb.v} & Clock and reset generator: 10\,ns period, reset asserted from time
  zero and released at 123\,ns, deliberately between clock edges. \\
\file{debug/sva/} & Assertion and coverage layer, attached with \sig{bind} so neither the RTL nor
  the testbench is modified. \\
\file{debug/hdl/tb\_axil\_master.vh} & Shared AXI4-Lite master driver tasks, so all five
  register-level benches drive the slave the same way and a change to the protocol handling
  happens once. \\
\file{debug/sim/rtl.f} & The RTL source list, ordered bottom-up by dependency level and shared by
  both flows; see Chapter~\ref{ch:buildorder}. \\
\file{debug/sim/asm.py} & Two-pass RV32I assembler with range-checked immediates and a symbol
  table the bench includes, so checks do not break when a label moves. \\
\bottomrule
\end{tabular}
\end{table}

\section{Method}

Five properties of the method matter more than the individual tests.

\paragraph{The program uses the features rather than poking signals.} Traps are
verified by taking them: handlers dispatch on \reg{mcause}, repair \reg{mepc} and
return, and execution continues.

\paragraph{Counts are exact, not approximate.} The direct-mode synchronous
handler counts every entry and must reach exactly 17; a register accumulates one
bit per cause and must equal \texttt{0x1FF}; and \reg{mhpmcounter6} independently
counts every trap entry of any kind and must equal exactly 26 (17 direct, 1
vectored, 8 interrupts). A missing, doubled or spurious trap cannot cancel out
against another.

\paragraph{Protocol legality is checked, not assumed.} The monitors run on every
link on every run, and the assertion layer re-states the same contract as
concurrent properties.

\paragraph{Configurations are swept.} The same suite runs under different bus
timing, so a result that depends on latency is exposed rather than hidden. Each
run echoes its parameters read back from the elaborated design, because a
silently ignored override looks exactly like a pass.

\paragraph{Coverage is measured and gated.} Required bins must be hit or the run
fails.

\paragraph{Negative properties get a witness, and a paired positive.} ``No
transaction was issued'' and ``the interrupt was never taken'' cannot be observed
directly, so the benches count what did happen --- AXI handshakes, acknowledge
pulses, handler entries --- and require the count to be zero. A count is
falsifiable; an absence of output is not. And every such check is followed by the
step that removes the suppressing condition and requires the thing to happen,
which is what separates masking from a dead path.

\paragraph{A test of a read-only register is only as good as its precondition.}
Writing zero to a register that already reads zero and then observing zero proves
nothing. Every read-only register is driven to a known non-zero value by real
hardware activity before the write is attempted.

\section{Regression results}

The following are the results of the regression as run on the documented
baseline.

\begin{table}[H]
\centering
\caption{ModelSim regression, twelve runs from one compile}
\label{tab:regress}
\begin{tabular}{@{}clL{0.36\linewidth}rl@{}}
\toprule
\textbf{Run} & \textbf{Bench} & \textbf{Configuration or scope} & \textbf{Checks} & \textbf{Result} \\
\midrule
1  & \file{tb\_cpu\_axi}    & Default latencies; the CPI check is active only here & 90 & Passed \\
2  & \file{tb\_cpu\_axi}    & High fixed AXI latencies (imem RL=2, dmem RL=3 / WL=2) & 90 & Passed \\
3  & \file{tb\_cpu\_axi}    & Random \sig{READY} backpressure, seeds 101 / 202 & 90 & Passed \\
4  & \file{tb\_cpu\_axi}    & Random \sig{READY} backpressure, seeds 777 / 888 & 90 & Passed \\
5  & \file{tb\_dual\_core}  & Two cores, shared memory, 2:1 arbiter & --- & Passed \\
6  & \file{tb\_pic}         & PIC feature bench & 139 & Passed \\
7  & \file{tb\_pic\_reset}  & PIC reset values, X-freedom, asynchronous reset & 350 & Passed \\
8  & \file{tb\_pic\_ro}     & PIC read-only registers and reserved bits & 478 & Passed \\
9  & \file{tb\_pic\_status} & PIC \reg{SRCx\_STATUS}, field by field & 424 & Passed \\
10 & \file{tb\_mtimer\_regs}& Machine-timer registers, reset and access rules & 224 & Passed \\
11 & \file{tb\_csr\_ro}     & CSR read-only / WARL / tied-off / absent & 282 & Passed \\
12 & \file{tb\_traps}       & One directed test per trap cause & 111 & Passed \\
\bottomrule
\end{tabular}
\end{table}

Runs 1--6 are the system and feature regression. Runs 7--12 are block-level
benches that isolate what a system run cannot: reset values, read-only
enforcement, individual status fields, and one trap cause at a time. The
procedures behind them are in Chapter~\ref{ch:verifplan}.

All twelve runs report \texttt{Errors: 0, Warnings: 0} at compile and elaboration.
The Verilator flow reports \texttt{ALL TESTS PASSED}, no assertion violations,
and no warnings on its default warning set.

\begin{table}[H]
\centering
\caption{Measured results}
\label{tab:measured}
\begin{tabular}{@{}lL{0.52\linewidth}@{}}
\toprule
\textbf{Metric} & \textbf{Value} \\
\midrule
Total self-checks in the regression & 2\,278, all passing \\
System bench checks    & 90, all passing \\
PIC feature bench      & 139 checks, all passing \\
PIC reset bench        & 350 checks, all passing \\
PIC read-only bench    & 478 checks, all passing \\
PIC status bench       & 424 checks, all passing \\
Machine-timer bench    & 224 checks, all passing \\
CSR access-rule bench  & 282 checks, all passing \\
Trap-cause bench       & 111 checks, all passing \\
Assertion violations   & 0 \\
Functional coverage    & 88 of 92 bins (95\,\%), gate passed \\
Synchronous trap count & Exactly 17 direct, 26 total \\
CPI, straight-line code, latency-1 memory & 1.00 (33 cycles for 33 instructions) \\
Misprediction cost     & 1 cycle \\
\sig{WFI} bus activity & At most one instruction fetch per sleep window \\
\bottomrule
\end{tabular}
\end{table}

The four uncovered bins are deliberate. Two are negative tests --- a channel that
is PIC-enabled but \reg{mie}-masked must \emph{never} be taken or acknowledged,
so its ``taken'' and ``acked'' bins must stay at zero. The other two are
AR-backpressure bins that the randomised regression configurations exercise
instead of the default run.

\begin{deviation}
\textbf{Not measured.} No synthesis was run, so there are no area, frequency or
power figures anywhere in this document. No riscv-arch-test suite or randomised
instruction stream has been run; the directed program is broad but hand-written.
Any statement about those would be unsupported and is therefore absent.
\end{deviation}

\section{Coverage of the PIC feature set}

The system bench runs the PIC in its default configuration and proves the
CPU--PIC contract. The feature bench covers the mechanisms individually.

\begin{table}[H]
\centering
\caption{PIC feature bench areas (\file{debug/hdl/tb\_pic.v})}
\label{tab:picverif}
\begin{tabular}{@{}L{40mm}L{0.52\linewidth}@{}}
\toprule
\textbf{Area} & \textbf{What is exercised} \\
\midrule
Reset defaults      & All 59 mapped registers read back against their documented reset values,
  twice: out of power-on reset and out of an asynchronous reset asserted off a clock edge from a
  fully configured, nested state. Nothing in the block leaves reset holding X, including the
  nesting-stack arrays, which are probed directly. \\
Read-only enforcement & Every read-only register and every reserved bit position: the write must
  be rejected with SLVERR \emph{and} the value must be unchanged, with the register first driven
  non-zero by hardware so the second half is not vacuous. \\
Status fields       & Each of the six \reg{SRCx\_STATUS} fields on its own and in both directions;
  the bump-escalation walk watched on every clock edge, because a register read is too slow to
  distinguish three correct steps from one jump. \\
Priority bands      & Inter-band, intra-band and index tie-break resolution. \\
Preemptive nesting  & A more urgent source preempting an active handler; depth tracking; \reg{NEST\_MAX} overflow raising \reg{OVF} without loss. \\
Band reconfiguration & A source re-banded while active keeps its claimed priority. \\
Spurious detection  & A source deasserting between offer and claim; \reg{SPURIOUS\_LOG} write-1-to-clear. \\
Deadline escalation & Escalation flipping the offer and preempting an active lower source; bump and multi modes. \\
Software triggers   & Keyed trigger accepted; unkeyed write ignored. \\
Edge sources        & Edge latched and consumed on claim; no re-fire. \\
AXI responses       & OKAY on mapped registers, SLVERR on unmapped offsets and read-only writes. \\
\bottomrule
\end{tabular}
\end{table}

\section{CPU verification areas}

\begin{table}[H]
\centering
\caption{Principal CPU verification areas and their evidence}
\label{tab:cpuverif}
\begin{tabular}{@{}L{30mm}L{0.40\linewidth}L{0.30\linewidth}@{}}
\toprule
\textbf{Area} & \textbf{Objective} & \textbf{Evidence} \\
\midrule
ISA coverage & Every instruction class executes correctly & Per-register result checks; instruction-class coverage bins \\
Forwarding   & S3$\to$S2 bypass into every consumer, zero stall cycles & Dependent chains; both-operand forwarding bin \\
Byte lanes   & \sig{WSTRB} and load extraction for all widths & Memory word checks plus sign/zero-extension register checks \\
Predictor    & Learning, not-taken default, alias correctness & Loop sum; three branches sharing an index with different tags \\
Traps        & Every synchronous cause, correct \reg{mepc} per cause & Cause bitmask \texttt{0x1FF}; exact count 17 \\
No side effects & An illegal store must not reach the bus & A word targeted by an illegal store stays zero \\
Alignment    & A misaligned access issues no transaction & Monitors observe no extra transaction \\
Bus errors   & SLVERR and DECERR on fetch, load and store & Accesses into unmapped space trap with causes 1, 5, 7 \\
CSR rules    & All six access forms plus negatives & Round-trips; illegal-access negatives; pure-read exception \\
Interrupts   & Sampled under MIE, at boundaries, held through stalls & Ack timestamp after the response beat; \reg{mepc} at a boundary \\
Counters     & 64-bit halves, M-mode writes, tied-off registers & Upper halves read 0; write/readback exact; tied-off registers read 0 and swallow writes without trapping \\
\sig{WFI}    & Sleep, silence, both wake paths & \reg{mepc} = \sig{wfi+4} on wake; fall-through with MIE = 0; bus quiet \\
RAS          & Returns from several call sites and a nested chain & Accumulated signature; return-prediction coverage bins \\
Multi-core   & Two cores, shared memory, flag handshake & Dual-core bench passes \\
\bottomrule
\end{tabular}
\end{table}

\section{Defects found by this environment}

Recording what the method caught is part of the evidence that it works.

\begin{table}[H]
\centering
\caption{Defects found, and by which pillar}
\label{tab:bugs}
\begin{tabular}{@{}L{22mm}L{0.56\linewidth}L{25mm}@{}}
\toprule
\textbf{Kind} & \textbf{Defect} & \textbf{Found by} \\
\midrule
RTL & Two AXI response-code flags (\sig{bresp\_err\_q}, \sig{rresp\_err\_q}) in the shared slave
  were unreset. \sig{BRESP} and \sig{RRESP} are combinational functions of them and are driven
  whether or not the corresponding \sig{VALID} is high, so both peripherals presented X on their
  response codes from time zero until their first transaction. A monitor cannot distinguish that
  from a real protocol violation. Fixed by resetting them, and by extending the reset policy to
  every flop in the delivery (Chapter~\ref{ch:reset}). & Reset bench, X-freedom checks \\
RTL & The interrupt controller's nesting-stack arrays were unreset, and are read by
  \reg{NEST\_STATUS.TOP\_ID} and \reg{ACTIVE\_VEC.ID}. Fixed by the same policy change. &
  Reset bench, direct probe of state with no register view \\
RTL & The bit positions the register map describes as reserved in \reg{SRCx\_CONFIG} ([15:8],
  [3]) and \reg{ESCALATION\_CFG} ([7:5], [3:2]) were stored and read back rather than hardwired
  to zero, so software could acquire a dependency on a bit a later revision might define. Fixed
  by masking on write. & Read-only bench, reserved-bit checks \\
RTL & \sig{ARVALID} asserted during reset. The issue logic is combinational and requested
  \sig{RESET\_PC} while \sig{rst\_n\_i} was still low. AXI forbids this, and a slave holding
  \sig{READY} high in reset would have accepted a phantom address. Fixed by gating with
  \sig{rst\_n\_i}. & Assertion layer, first run \\
RTL & Illegal load/store encodings issued their AXI transaction before trapping, so an illegal
  store could write memory. Fixed by gating the LSU request with the decoder's \sig{illegal}. &
  Scoreboard check \\
Testware & The assembler silently truncated an out-of-range label immediate. Now range-checked. &
  Directed test failure \\
Testware & The memory model could not perform latency-1 reads, so the CPI check was measuring the
  model rather than the CPU. & CPI check review \\
Testware & An X leaked through the testbench address decoder. & Random-backpressure runs \\
Infrastructure & \texttt{-g} parameter overrides were silently swallowed by \texttt{vopt}, so one
  regression configuration ran at defaults while reporting green. Switched to \texttt{-G} plus a
  per-run parameter echo. & Configuration sweep \\
Test plan & The first coverage measurement (59 of 85 bins) exposed real holes: \sig{AUIPC} never
  committed, four branch forms never executed, \sig{CSRRCI} never executed, both-operand forwarding
  never hit, and five PIC channels never taken. All closed. & Functional coverage \\
Integration & A source enabled in the PIC but masked in \reg{mie} starves every less urgent source.
  Not an RTL defect; documented as an integration rule and respected by the test program. &
  Channel sweep \\
\bottomrule
\end{tabular}
\end{table}

A note on the first row, because it is the clearest argument for the assertion
layer existing at all: the procedural monitor could not have found it. Monitors
arm after reset, and the violation occurs during reset.
